\documentclass[a4paper,11pt]{article}
\usepackage[margin=1cm, bottom=1.5cm]{geometry}
\usepackage[utf8]{inputenc}
\usepackage[T1]{fontenc}
\usepackage{helvet}
\renewcommand{\familydefault}{\sfdefault}
\usepackage{xcolor}
\usepackage{tikz}
\usepackage{multicol}
\usepackage{graphicx}
\usepackage[hidelinks]{hyperref}
\usepackage{enumitem}

% Define custom colors from design tokens
\definecolor{sage}{RGB}{132, 169, 140}
\definecolor{lavender}{RGB}{216, 216, 246}
\definecolor{coral}{RGB}{244, 151, 142}
\definecolor{canvas}{RGB}{253, 252, 248}
\definecolor{darktext}{RGB}{50, 60, 55}

% Set background color
\pagecolor{canvas}

\pagestyle{empty}

\begin{document}

% --- Header with Logo and Title ---
\noindent
\begin{tikzpicture}
    % Background box for header
    \fill[sage!15, rounded corners=1cm] (0,0) rectangle (\textwidth, 4.2);
    
    % Logo (Left)
    \node[anchor=west] at (0.5, 2.1) {\includegraphics[height=3.2cm]{../public/BIOMAP_logo.png}};
    
    % Date Box (Right) - Square, same height as logo (3.2cm)
    \fill[coral, rounded corners=0.3cm] (\textwidth-3.7cm, 0.5) rectangle (\textwidth-0.5cm, 3.7);
    \node[align=center] at (\textwidth-2.1cm, 2.1) {\color{white}\textbf{\large AUG}\\[0.2em]\textbf{\fontsize{30}{36}\selectfont 22}};
    
    % Title Block (Middle)
    % x=4.0 reduced gap.
    \node[anchor=west, align=left] at (4.0, 2.1) {
        {\fontsize{36}{42}\selectfont \textbf{\color{sage!60!black}BIOMAP 2026}}\\[0.2em]
        {\fontsize{13.5}{16}\selectfont \textbf{\color{coral!90!black}BIO-inspired Methods for Pattern Recognition}}\\[1em]
        {\fontsize{14}{16}\selectfont \color{darktext} \textbf{Lyon, France}}
    };
\end{tikzpicture}

\vspace{0.3cm}

% --- Motivation ---
\begin{center}
    \parbox{\textwidth}{
% \centering removed for left alignment
        \small \color{darktext}
        
        \textbf{Pattern recognition systems in nature show strong abilities through evolutionary adaptation and self-organisation, from the visual cortex's layered feature learning to the group behaviour of swarm organisms.}
        \vspace{0.2cm}
        
        This workshop sets bio-inspired intelligence as a core approach for future pattern recognition systems, gathering views from evolutionary computation, swarm intelligence, artificial life, and neuromorphic computing. It covers a wide range of bio-inspired methods and their combined use. The workshop examines the key questions where biology and computation intersect: How can evolutionary ideas aid in the automatic discovery of pattern recognition structures? What does swarm intelligence show about shared visual processing? How do self-organising ideas allow strong pattern recognition in changing settings? It aims to start new ideas that cross standard field limits.

        \vspace{0.2cm}
        The workshop links to real-world problems in pattern recognition. For systems operating in unknown settings where trained models are ineffective, evolutionary adaptation and self-organisation provide avenues for learning and adjustment as they run. In medical imaging, where marked data is rare, bio-inspired methods allow learning from a few examples. For IoT applications, swarm intelligence provides scalable solutions. The workshop will demonstrate how these methods address accuracy, efficiency, robustness, and adaptability simultaneously.
        
        \vspace{0.2cm}
        The workshop also ties to themes of sustainable and trustworthy AI. Biological systems achieve high efficiency through evolution, providing insights for energy-efficient recognition. By establishing bio-inspired intelligence as a unified framework, this workshop will accelerate the development of robust, adaptive, efficient, transparent, and human-aligned pattern recognition systems.
    }
\end{center}

\vspace{0.2cm}

% --- Important Dates  ---
\begin{center}
    {\Large \textbf{\color{sage!60!black}Important Dates}}
\end{center}
\vspace{0.2cm}
\begin{tikzpicture}
    \node[fill=lavender!40, rounded corners=0.5cm, inner sep=0.4cm, minimum width=\textwidth] {
        \begin{tabular*}{\dimexpr\textwidth-1cm\relax}{@{\extracolsep{\fill}}ccc}
            \textbf{\color{sage!60!black} Paper Submission} & \textbf{\color{sage!60!black} Notification} & \textbf{\color{sage!60!black} Camera-Ready} \\
            {\large \color{coral!90!black} \textbf{May 1, 2026}} & {\large June 10, 2026} & {\large June 18, 2026}
        \end{tabular*}
    };
\end{tikzpicture}

\vspace{0.5cm}

% --- Topics ---
\begin{center}
    {\Large \textbf{\color{sage!60!black}Workshop Topics}}
\end{center}
\vspace{0.2cm}
\begin{multicols}{2}
    \begin{itemize}[label={\color{coral}\textbullet}, itemsep=0.0em, leftmargin=*]
        \small
        \item Evolutionary neural architecture discovery
        \item Swarm intelligence for distributed patterns
        \item Self-organising and adaptive vision
        \item Artificial life and emergent patterns
        \item Evolutionary multi-objective optimisation
        \item Bio-inspired learning paradigms
        \item Hybrid evolutionary-gradient methods
        \item Neuromorphic evolutionary systems
        \item Collective intelligence for large-scale analysis
        \item Dynamic and online evolutionary adaptation
        \item Evolutionary robotics and embodied vision
        \item Quantum-inspired evolutionary recognition
        \item Evolutionary explainability
        \item Bio-inspired hardware-co-evolution
        \item Applications (medical, remote sensing, etc.)
    \end{itemize}
\end{multicols}

\vspace{0.2cm}

% --- Organizers (Reordered: Fontanella, Nardone, Vanneschi) ---
\begin{center}
    {\Large \textbf{\color{sage!60!black}Workshop Chairs}}
\end{center}
\vspace{0.2cm}

\noindent
\begin{minipage}{0.33\textwidth}
    \textbf{Francesco Fontanella}\\
    \small University of Cassino\\and Southern Lazio
\end{minipage}%
\begin{minipage}{0.33\textwidth}
    \centering
    \textbf{Emanuele Nardone}\\
    \small University of Cassino\\and Southern Lazio
\end{minipage}%
\begin{minipage}{0.33\textwidth}
    \raggedleft
    \textbf{Leonardo Vanneschi}\\
    \small NOVA IMS,\\Universidade Nova de Lisboa
\end{minipage}

\vspace{0.3cm}

% --- Footer ---
\begin{center}
    \begin{tikzpicture}
        \fill[sage!15, rounded corners=0.2cm] (0,0) rectangle (10, 0.9);
        \node at (5, 0.45) {
            \large \textbf{\color{sage!60!black} \href{https://narden91.github.io/biomap}{Visit the BIOMAP Website}}
        };
    \end{tikzpicture}
    
    \vspace{0.1cm}
    {\tiny \color{darktext} Accepted papers will be included in the ICPR 2026 Conference Proceedings.}
\end{center}

\end{document}